\documentclass[11pt,a4paper]{article}
\usepackage[spanish,es-noquoting]{babel}
\usepackage[top=2.5cm, bottom=2.5cm, left=2.5cm, right=2.5cm]{geometry}
\usepackage{fontspec}
\usepackage{amsmath}
\usepackage[version=4]{mhchem}
\usepackage{enumitem}
\usepackage{microtype}

\setmainfont{Calibri}
\linespread{1.15}

\begin{document}
\pagestyle{empty}

\noindent\textbf{1.} En un matraz se hacen reaccionar $49\text{ g}$ de ácido sulfúrico (\ce{H2SO4}) con $32\text{ g}$ de hidróxido de sodio (\ce{NaOH}, sosa cáustica) según la siguiente reacción de neutralización (sin ajustar):
\[
\ce{H2SO4 + NaOH -> Na2SO4 + H2O}
\]

\begin{enumerate}[label=\textbf{\alph*)}, leftmargin=2em, itemsep=6pt]
    \item Escribe y ajusta la ecuación química correspondiente a la reacción de neutralización.
    \item Determina de forma justificada cuál de los dos reactivos actúa como reactivo limitante y cuál en exceso.
    \item Calcula la masa de sulfato de sodio (\ce{Na2SO4}) que se obtendrá al completarse la reacción.
    \item Determina los moles y los gramos del reactivo en exceso que quedan sin reaccionar.
    \item Calcula la masa de agua obtenida y comprueba la ley de conservación de la masa (Lavoisier).
\end{enumerate}

\vspace{0.4cm}
\noindent\textbf{Datos:} Masas atómicas ($\text{u}$): $\mathrm{H} = 1$; $\mathrm{O} = 16$; $\mathrm{Na} = 23$; $\mathrm{S} = 32$.

\end{document}
